\documentclass{../base/canon}

%% canon_variables.tex — Valores por defecto de metadatos institucionales
%%
%% Incluir ANTES del \begin{document} para sobreescribir los defaults de la clase.
%% El template puede sobreescribir cualquier valor después de este \input.
%%
%% Uso en template:
%%   %% canon_variables.tex — Valores por defecto de metadatos institucionales
%%
%% Incluir ANTES del \begin{document} para sobreescribir los defaults de la clase.
%% El template puede sobreescribir cualquier valor después de este \input.
%%
%% Uso en template:
%%   %% canon_variables.tex — Valores por defecto de metadatos institucionales
%%
%% Incluir ANTES del \begin{document} para sobreescribir los defaults de la clase.
%% El template puede sobreescribir cualquier valor después de este \input.
%%
%% Uso en template:
%%   \input{../base/canon_variables}
%%   \SetDocTitle{Mi Informe Específico}   % sobreescribe el default

\SetDocLogo{../assets/logo.pdf}
\SetDocUnit{Tecnología \& Sistemas}
\SetContactUnit{Mesa de soporte}
\SetContactMembers{%
  Equipo CoreX & Soporte & soporte@empresa.com%
}
\SetContactNote{Para soporte técnico o consultas sobre este documento.}

%%   \SetDocTitle{Mi Informe Específico}   % sobreescribe el default

\SetDocLogo{../assets/logo.pdf}
\SetDocUnit{Tecnología \& Sistemas}
\SetContactUnit{Mesa de soporte}
\SetContactMembers{%
  Equipo CoreX & Soporte & soporte@empresa.com%
}
\SetContactNote{Para soporte técnico o consultas sobre este documento.}

%%   \SetDocTitle{Mi Informe Específico}   % sobreescribe el default

\SetDocLogo{../assets/logo.pdf}
\SetDocUnit{Tecnología \& Sistemas}
\SetContactUnit{Mesa de soporte}
\SetContactMembers{%
  Equipo CoreX & Soporte & soporte@empresa.com%
}
\SetContactNote{Para soporte técnico o consultas sobre este documento.}

\SetDocType{FICHA RESUMEN}
\SetDocTitle{Título de la Ficha Resumen}
\SetDocCode{FIC-XXX-000}
\SetDocArea{Área}
\SetDocAuthor{Autor}
\SetDocDate{\today}

\begin{document}

\section{Indicadores clave}

\begin{center}
\begin{tabular}{ccc}
  \FichaKPI{Indicador 1}{0}{período} &
  \FichaKPI{Indicador 2}{0\,\%}{meta: 0\,\%} &
  \FichaKPI{Indicador 3}{0}{límite: 0} \\[1.4em]
  \FichaKPI{Indicador 4}{0}{referencia} &
  \FichaKPI{Indicador 5}{0\,\%}{meta: 0\,\%} &
  \FichaKPI{Indicador 6}{0}{unidad} \\
\end{tabular}
\end{center}

\section{Nota de contexto}

\begin{ParrafoEjecutivo}
Una o dos oraciones de contexto para quien lee esta ficha sin haber visto el
informe completo. Indicar el período y cualquier condición relevante.
\end{ParrafoEjecutivo}

\ObservationBox[Observación]{
  Observación o alerta más relevante del período resumido en esta ficha.
}

\sectionrule

\begin{center}
  {\footnotesize\color{CanonMuted}
    Ficha generada el \@docdate{} — para el informe completo consultar \texttt{\@doccode}}
\end{center}

\end{document}
